% Draft for Section 4 (Evolutionary Process, Variation Operators, and Population Evidence)
\section{Evolutionary Process, Variation Operators, and Population Evidence}
\label{sec:evolution}

\subsection{Implemented Variation Operators: Block-Preserving Crossover and Six-Mode Mutation}
Rather than treating evolutionary search as an opaque black-box optimizer, Jarambit records the exact variation operators applied across its island populations:
\begin{itemize}[leftmargin=*,nosep]
\item \textbf{Two-parent crossover:} To prevent non-physical saw-tooth discontinuities along the closed silhouette spline, the 40 normal control points (\code{cp00}--\code{cp39}, indices $0..39$) are inherited as a single contiguous block from parent~$b$ with probability $0.5$ and otherwise retained from parent~$a$. The remaining 88 numeric feature coordinates (indices $40..127$) and the categorical alloy grade (\code{material}) are inherited independently with probability $0.5$ per coordinate, followed immediately by deterministic clamping.
\item \textbf{Six-mode structured mutation:} Each mutation call draws an integer mode $\text{op}\in\{0,1,2,3,4,5\}$ uniformly at random with step scale $\sigma=0.15$ (coarse) or $\sigma=0.04$ (fine):
  \begin{enumerate}[leftmargin=*,nosep]
  \item \emph{Mode~0 (Sparse per-gene Gaussian):} Perturbs each coordinate $i\in[0,127]$ independently with probability $0.15$ by $\mathcal{N}(0,1)\cdot(\text{hi}_i-\text{lo}_i)\sigma$.
  \item \emph{Mode~1 (Smooth silhouette bump):} Applies a periodic Gaussian normal displacement bump $A\exp(-d^2/(2w^2))$ across all 40 control points around a random station $c\in[0,39]$ with width $w\in[1.5,5.5]$.
  \item \emph{Mode~2 (Thickness, taper, and relief operator):} Perturbs \code{t\_stock}, \code{t\_tip}, \code{taper\_x0}, \code{relief\_d}, and \code{relief\_w}, and---when \code{p1\_on} $\ge 0.5$---also perturbs \code{rail\_w}, \code{pocket\_d}, and one interior spline knot \code{k0}--\code{k4}.
  \item \emph{Mode~3 (Dimple-lattice operator):} Deactivates the lattice (\code{lattice\_d = 0}) with probability $0.15$ or perturbs lattice depth, pitch, fill fraction, angle, and handle/blade extents ($68..73$).
  \item \emph{Mode~4 (Fuller, flute, texture, and comet operator):} Perturbs blade fuller and handle flute geometry ($74..78$), surface texture mode and pitch/angle ($79..82$), and comet groove width/bow ($52..54$).
  \item \emph{Mode~5 (Internal pocketing or continuous-rail operator):} When \code{p1\_on} $\ge 0.5$, mutates the continuous perimeter rail and knot coordinates ($118..127$) with probability $0.7$; otherwise selects one of four sub-operators targeting triangular isogrid ($90..98$), pierced windows ($99..104$), ring flute ($105..109$), or angled micro-perforations ($110..116$).
  \end{enumerate}
  In coarse mode, the alloy grade flips between \code{Ti-6Al-4V} and \code{Ti-3Al-2.5V} with probability $0.03$ before clamping.
\end{itemize}

\subsection{Four Distinct Population Collections}
Scientifically credible reporting requires distinguishing four collections that share candidates but have different denominators and selection biases:
\begin{enumerate}[leftmargin=*,nosep]
\item \textbf{Birth collection:} Every proposed candidate evaluated or rejected at birth:
  \begin{itemize}[leftmargin=*,nosep]
  \item \emph{Cohort~3 (10-Run 800-Gen Population, $1{,}954{,}560$ Evaluations):} \text{Run~2} through \text{Run~11} each evaluate $195{,}456$ births across generations $0$--$800$ ($801$ distinct generations per run), totaling $1{,}954{,}560$ complete 146-column birth records ($754{,}236$ search-legal).
  \item \emph{Cohort~1 (Historical 17-Gen, \text{Run~1}):} $1{,}984$ births across generations $0$--$16$ ($1{,}455$ with completed plate finite-element solves, $529$ not evaluated by finite elements, and $446$ search-legal).
  \item \emph{Cohort~2 (Intermediate Screening):} $3{,}649$ multi-stage birth proposals across 24 screening runs, with $2{,}259$ joined to Stage-1 screening observations and $1{,}390$ preserved in an explicit unscreened observation band.
  \end{itemize}
\item \textbf{Island-membership (survivor) collection:} Candidates occupying active island slots at generation end, including elite retention and periodic ring migration. In Cohort~1 (\text{Run~1}), $2{,}176$ slot occurrences are recorded across Islands~1--4 ($4\times 32$ slots over 17 generations), comprising $1{,}984$ initial placements, $192$ elite reappearances, and $12$ ring-migration events ($4$ events each at generations 4, 8, and 12).
\item \textbf{Quality-diversity archive:} Representatives retained per behavioral niche cell across each run.
\item \textbf{Independent Stage-2 / strict assessment collection:} The subset of top-ranked or editorially probed candidates admitted to full 26-gate CAD, plate, and quadratic solid verification.
\end{enumerate}

\begin{figure*}[t]
\centering\includegraphics[width=\textwidth]{population.png}
\caption{Cohort~1 (\text{Run~1}) historical birth and island-membership collections. All $1{,}984$ retained births across generations 0--16 ($1{,}455$ finite-element-complete blue points, $529$ unevaluated amber crosses, $446$ search-legal births) are plotted on an $\operatorname{asinh}(\text{fit})$ axis that preserves the negative legality-penalty tail without clipping. Recorded four-island slot memberships ($2{,}176$ occurrences across Islands~1--4) preserve $192$ elite reappearances and $12$ ring-migration events at generations 4, 8, and 12.}
\label{fig:population}
\end{figure*}

\subsection{Recorded Multi-Hop Inheritance and Coordinate-Level Genome Diffs}
Figure~\ref{fig:lineage} traces the actual three-hop ancestry of Individual~1095 (Cohort~1, \text{Run~1}, generation~9, highest stored birth fitness $\text{fit}=1.15727$). Rather than a simple tree of independent ancestors, the recorded parent pointers form a directed acyclic graph with ancestor reuse:
\begin{itemize}[leftmargin=*,nosep]
\item Individual~174 (generation~1, $\text{fit}=1.14979$, archive crossover-mutation from parents 57 and 61) mutates to produce both Individual~245 (generation~2, $\text{fit}=1.15388$, 24 changed \code{f64} coordinates, $L_2$ gene distance $0.604263$) and Individual~299 (generation~2, $\text{fit}=1.14855$, 1 changed coordinate \code{t\_stock}, $\Delta = +0.032993\,\mathrm{mm}$).
\item Individual~245 mutates into Individual~413 (generation~3, $\text{fit}=1.15296$, 5 changed coordinates, $L_2=0.903516$), while Individual~299 mutates into Individual~392 (generation~3, $\text{fit}=1.12512$, 6 changed coordinates, $L_2=4.948287$).
\item Individual~526 (generation~4, $\text{fit}=1.15563$) arises via crossover-mutation from parents 413 and 392 (differing in 30 coordinates from Individual~413 and 22 coordinates from Individual~392).
\item Individual~1095 (generation~9, $\text{fit}=1.15727$) arises via crossover-mutation from parents 413 and 526, reusing ancestor Individual~413 across two generations and differing in 30 coordinates from Individual~413 ($L_2=13.720172$) and 5 coordinates from Individual~526 ($L_2=0.946269$).
\end{itemize}
Crucially, when Individual~1095 was admitted to strict multi-model verification, it failed the regional failure-sequence requirement (\code{cert\_pass\_all=false}), demonstrating why high search-stage fitness and recorded lineage do not establish engineering qualification.

\begin{figure*}[t]
\centering\includegraphics[width=\textwidth]{lineage.png}
\caption{Recorded genealogy and coordinate-level inheritance differences around Individual~1095 (Cohort~1, \text{Run~1}, generation~9, $\text{fit}=1.15727$). Directed edges preserve ancestor reuse (Individual~413 and Individual~174) and exact changed-coordinate counts across all 128 numeric genes and the alloy column.}
\label{fig:lineage}
\end{figure*}
